\subsection{Die vollständige Neutrinomatrix besitzt ein norm- und phasentreues Quellenwörterbuch}\label{sec:neutrinodictionary}
Die vorhandenen Neutrinoherleitungen liefern ganze komplexe Matrizen.
Ihr Quellenanschluss darf sich deshalb nicht auf Massen, Spur oder Norm
beschränken. Der bereits konstruierte symmetrische Sechserraum hat genau
die dafür benötigte bilineare Struktur. Für eine komplex-symmetrische
Koeffizientenmatrix $A$ sei
\begin{equation}
 \Gamma(A)=\sum_a A_{aa}B_{aa}+\sum_{a<b}A_{ab}B_{ab}.
\end{equation}
Die 36 ursprünglichen affinen Paarungen ergeben
\begin{equation}\label{eq:pairisometry}
 \langle\Gamma(A),\Gamma(B)\rangle=\tr(A^\dagger B).
\end{equation}
Dabei haben diagonale Paare Normquadrat eins und außerdiagonale Paare
Normquadrat zwei. Die rechte Seite ist deshalb genau das volle
Frobenius-Skalarprodukt, einschließlich der komplexen Kreuzterme.

Auch das tatsächliche Auslesen ist berechnet. Mit $C(v)=\sum_a v_a C_a$ gilt
\begin{equation}\label{eq:pairtomography}
 \langle\Gamma(A)^\dagger(\infty)C(v)(z_1)C(v)(z_2)\rangle
 =(z_1-z_2)^2\,v^{\mathsf T}A^\dagger v.
\end{equation}
Die sechs kohärenten Einsetzungen
$e_1,e_2,e_3,e_1+e_2,e_1+e_3,e_2+e_3$ rekonstruieren daher jedes $A$.
Die gemessene Größe ist hierbei eine \emph{komplexe Amplitude mit Phasenreferenz}.
Sechs Betragsquadrate reichen nicht: Sie unterscheiden im Allgemeinen eine
Matrix und ihre komplex konjugierte Matrix nicht.

Der Familienbasiswechsel wird aus den vorhandenen Generatoren konstruiert.
Setze $F=(e_1,Me_1,M^2e_1)$ und $S=F^\dagger J_C$, wobei $J_C$ der zuvor
an den Originalströmen geprüfte Intertwiner ist. $S$ ist unitär und verbindet
den Quellenvektor $v$ mit der ursprünglichen zyklischen Darstellung $w=Sv$.
Die einzelne C-Wirkung behält ihr Minuszeichen gegenüber dem Familienumlauf;
im symmetrischen Paar hebt es sich auf. Ein kovarianter physischer
Massentensor $K$ zieht sich dagegen als $K_C=S^{\mathsf T}KS$ zurück.
Der richtige Zustandsanschluss ist deshalb
\begin{equation}\label{eq:massformencoding}
 \mathcal E(K)=\Gamma\!\left(\overline{S^{\mathsf T}KS}\right),\qquad
 \frac{\langle\mathcal E(K)^\dagger(\infty)C(v)(z_1)C(v)(z_2)\rangle}
 {(z_1-z_2)^2}=(Sv)^{\mathsf T}K(Sv).
\end{equation}
Die Konjugation ist wesentlich: $\Gamma$ ist komplexlinear auf
Zustandskoeffizienten, $\mathcal E$ dagegen antilinear auf Massentensoren.
Es gilt $\langle\mathcal E(K),\mathcal E(L)\rangle=\tr(L^\dagger K)$.
Ein untypisierter Ähnlichkeitstransform würde diese physische bilineare
Antwort nicht erhalten.

\paragraph{Schwere und leichte Neutrinos bleiben verschiedene physische Felder.}
Der konkrete D5-Singulettpaarkanal besitzt $X=+10$ und gehört zum
$\nu^c\nu^c$- beziehungsweise konjugierten schweren Majorana-Anschluss.
Sein direkter Familientensor ist die ursprüngliche schwere Form
\[
 M_R=M_{\rm scal}\,\diag(\varepsilon,2\varepsilon,3),
 \qquad \varepsilon=\phi_0^2/(2g_{\rm car}),
\]
in der erklärten Q$_+$-Eigenbasis. Auch sie wird durch
$\mathcal E(M_R)$ samt allen sechs Rückleseamplituden dargestellt.
Für den dimensionslosen Tensor $D_M=M_R/M_3$ ergibt sich ausdrücklich
\[
 D_M=\diag(\varepsilon/3,2\varepsilon/3,1),\qquad
 \norm{\mathcal E(D_M)}^2=1+5\varepsilon^2/9.
\]
Das ist die direkte schwere Quellenübersetzung. Die anschließende bereits
verwendete normierte Seesaw-Rechnung lautet
\[
 K_\nu=-\widehat Y^{\mathsf T}D_M^{-1}\widehat Y,\qquad
 \widehat Y=i\sqrt{D_M}\,R_{\rm CI}\sqrt{m_{\rm niedrig}}\,U^\dagger.
\]
Hier hat $\widehat Y$ in der erklärten gemeinsamen RG-Normierung die Einheit
$\sqrt{\rm eV}$, während $K_\nu$ in eV steht. $\widehat Y$ ist daher nicht
mit dem dimensionslosen physischen Yukawa-Operator gleichzusetzen.
Die vorhandene Wahl der gemischten $\varepsilon$-Einsetzung bleibt dabei
ein ausgewiesener Ansatz; ihre Übereinstimmung mit dem physisch markierten
Quellenrahmen wird nicht allein durch einen Basiswechsel bewiesen.
Der leichte Massentensor entsteht anschließend über die vorhandene
Seesaw-Regel aus $M_R$ und $Y_\nu$.
Ein leichter $3\times3$-Tensor kann mathematisch in denselben Sechserraum
kodiert werden. Daraus folgt aber keine direkte Identifikation eines
linken Standardmodell-Neutrinopaares mit dem geladenen $X=10$-Quellenoperator.
Diese Feldtypgrenze wird in der folgenden Phasenprobe erhalten.

\paragraph{Die bereits vorhandenen gemeinsamen Neutrinobedingungen bleiben erhalten.}
Das Programm verwendet den ursprünglichen Operator $R_{\rm CI}(z)$ und
die eingefrorenen Eingaben von \file{nu_scalaron_v3.yaml}.
Die bestehende gleichzeitige Bedingung aus größtem Dirac-Singularwert und
$\widetilde m_1=m_3/10$ aus Abschnitt~\ref{sec:neutrinojoint} ergibt erneut
\[
 m_3=0{,}0435180861582\ {\rm eV},\qquad |z|=0{,}477382171162.
\]
Dies ist die bereits dokumentierte bedingte Konsistenzlösung, keine neue
Massenprognose. Insbesondere bleiben die Kalibrierung von $m_2$, der
Washout-Ansatz, die gewählte CI-Symmetrieklasse und der erklärte gemeinsame
Skalenfaktor erhalten; vollständige Schwellen-RG- und Flavor-Boltzmann-Dynamik
werden dadurch nicht zusätzlich geprüft.

\paragraph{Die verbleibende Phase wird nun zu einer konkreten Quellenantwort.}
Im ausdrücklich phasenoffenen Zweig mit $m_1=0$ verbleibt eine physische
Majorana-Phase. Mit dem ursprünglichen $U_0$ lautet diese Familie
\begin{equation}
 K_\nu(\gamma)=\overline{U_0}\,\diag(0,m_2e^{-i\gamma},m_3)\,U_0^\dagger.
\end{equation}
Die unveränderten leichten Massen, PMNS-Winkel und Dirac-CP-Invariante,
$YY^\dagger$, die ungeflavorten schweren CP-Invarianten sowie
$\norm{\mathcal E(K_\nu)}^2=m_2^2+m_3^2$ wählen $\gamma$ nicht aus.
Mit $v_j=S^\dagger u_j$, wobei $u_j$ die Spalten von $U_0$ sind,
liefern die ursprünglichen Paaramplituden dagegen
\begin{equation}\label{eq:majoranaphaseresponse}
 \frac{\mathcal A_2}{\mathcal A_3}
 =\frac{m_2}{m_3}e^{-i\gamma}.
\end{equation}
Die volle Quelle unterscheidet somit genau die relative Phase, die ihre
neutralen Normantworten ausblenden. Dieser Befund gilt für die im YAML
ausdrücklich als unbeschränkte Majorana-Phasen geführte Klasse.
Wird die gedruckte vollständige komplexe Yukawa-Matrix als unveränderlicher
Ansatz übernommen, ist $\gamma=0$ bereits durch diesen Ansatz festgelegt.
Eine Variation von $\gamma$ erhält dann nicht alle gedruckten Matrixeinträge.
Da $YY^\dagger$ unverändert bleibt, behebt die Phase auch nicht das zuvor
gefundene Defizit der ungeflavorten gemeinsamen Baryogeneserechnung.

\paragraph{Welche Symmetrie die Matrix tatsächlich auswählen darf.}
Die beiden ursprünglichen markierten Generatoren liefern einen zusätzlichen
exakten Test. In der zyklischen Darstellung seien sie $P$ und $Q$.
Sie erfüllen $P^3=Q^4=(PQ)^4=I$.
Für eine gemeinsame Charakterlinie des symmetrischen Paarraums folgen daraus
$\lambda^3=\mu^4=(\lambda\mu)^4=1$, also $\lambda=1$.
Eine $P$-invariante symmetrische Matrix besitzt die Form
\[
 A=\begin{pmatrix}a&b&b\\b&a&b\\b&b&a\end{pmatrix}.
\]
Die tatsächliche $Q$-Wirkung ersetzt sie durch
$\left(\begin{smallmatrix}a&b&-b\\b&a&-b\\-b&-b&a\end{smallmatrix}\right)$.
Proportionalität zu $A$ erzwingt $b=0$ und bei $A\ne0$ den Charakter $\mu=1$.
Die einzige gemeinsame Linie ist somit $A=aI$; in der ursprünglichen
C-Basis lautet sie
\[
 A_C=a\begin{pmatrix}1&0&0\\0&0&1\\0&1&0\end{pmatrix}.
\]
Ihr Rang ist drei, ihre drei Takagiwerte sind gleich.
Der vorhandene leichte Massentensor besitzt hingegen Rang zwei und
Takagiwerte $(0,m_2,m_3)$ mit $m_2\ne m_3$, für jede Majorana-Phase.
Eine zusätzlich geforderte vollständige Invarianz würde diesen vorhandenen
Tensor daher ausschließen.

Diese zusätzliche Invarianz steht nicht im ursprünglichen Neutrinovertrag:
NU.CI.01 verlangt die Zyklus- und Realitätsbedingungen an $R_{\rm CI}$;
die gedruckte schwere Matrix ist eine Generationenraum-/Q$_+$-Einsetzung.
Die passende weiterzuprüfende Quellenregel ist deshalb eine
\emph{kovariante} Abhängigkeit von den markierten Defekten $\mathcal D$,
\[
 K(g\mathcal D)=\rho(g)^{-\mathsf T}K(\mathcal D)\rho(g)^{-1},
\]
sofern die Quellenwirkung diese Kovarianz realisiert. Sie verlangt keinen
konstanten Fixvektor. Die Markierung kann die Familienrichtung unterscheiden;
die physische Higgs-/Masseneinsetzung muss ihre geladene Kopplung liefern.

\paragraph{Der jetzt genau bestimmte Auswahlgegenstand.}
Alle sechs Paare tragen $X=10$, während die neutralen Felder und die
Familien-Defekte diese D5-Ladung nicht kompensieren. Im ursprünglichen
$X$-invarianten Vakuum gilt folglich
\[
 \langle\mathcal E(K)\rangle=0.
\]
Eine Matrix als geladene Einsetzung zu berechnen ist damit etwas anderes
als ihre Auswahl durch einen physisch realisierten Hintergrund.
Der verschwindende Einpunktwert schließt einen eichinvarianten
Higgs-/Majorana-Mechanismus nicht aus. Der physische Anschluss kann einen
zugehörigen Higgsinsertionskanal oder einen insgesamt neutralen höheren
Korrelator verwenden; dann muss genau dessen ursprüngliche Abbildung
angegeben und ausgewertet werden.
Gesucht ist der tatsächlich erzeugte Koeffizient in \emph{diesem}
symmetrischen Sechserkanal, mit seiner Skala und Orientierung sowie dem
zugehörigen ladungskompensierenden Anschluss. Die vorhandene vollständige
Neutrinoherleitung legt fest, welchen Tensor dieser Anschluss reproduzieren muss.
\file{tfpt_explorer/source_neutrino_dictionary.py} berechnet nun seine
Isometrie, die ursprüngliche Generatorabbildung und die vollständige
phasensensitive Rücklesung.

\picturetext{Die Massentabelle lässt sich jetzt verlustfrei in die gemeinsame
Quelle übersetzen und wieder auslesen. Auch die Informationen, die beim
bloßen Wiegen der Zustände verloren gehen, bleiben erhalten. Der nächste
Herkunftsschritt muss erklären, warum der ursprüngliche Prozess genau
diese Tabelle in genau diesen Quellenzustand schreibt.}

\subsection{Ein gemeinsamer Transport für primitive Windung, Massentensor und neutrale Antwort}\label{sec:masssourceconnection}
Die Fortsetzung vom 29.~September arbeitet unmittelbar am vorher benannten
Anschluss. Die ursprüngliche Fuchsverbindung bestimmt nicht nur die
Familienendphasen. Ihr tatsächlicher Spin-Lift induziert auf den sechs
symmetrischen Paaren eine Verbindung, aus der sich die \emph{ganze}
ursprüngliche Verbindung wieder zurückgewinnen lässt. Auch die primitive
Determinantenwindung und die bereits berechnete neutrale RHP-Antwort lassen
sich an diesem selben Transport ablesen.

\paragraph{Die ursprüngliche Dreierwirkung muss mit ihrem Spin-Lift eingehen.}
Setze $d=\det\Psi$, $\Psi'=A\Psi$ und $t=\tr A$.
Die bereits an den tatsächlichen C-Strömen geprüfte Basisabbildung ist
$J=HP=\operatorname{diag}(1,-1,1)$, wobei $P$ die native Reihenfolge
$(23,13,12)$ in die Außenpotenzreihenfolge $(12,13,23)$ übersetzt und
$H$ die dreidimensionale Hodge-Abbildung ist. Die Identität
\[
 H(\Lambda^2\Psi)H^{-1}=d\,\Psi^{-\mathsf T}
\]
ergibt auf dem behaltenen C-Dreier den komplexifizierten internen Spintransport
\begin{equation}\label{eq:massspinconnection}
 S_C=d^{-1/2}P^{\mathsf T}(\Lambda^2\Psi)P
     =d^{1/2}J^{\mathsf T}\Psi^{-\mathsf T}J,
 \qquad
 B:=S_C'S_C^{-1}=\frac t2 I-J^{\mathsf T}A^{\mathsf T}J.
\end{equation}
Die Quadratwurzel wird entlang des ursprünglichen Wegs fortgesetzt.
Am unitären Ende ist das genau der schon nachgewiesene gewundene Spin-Lift;
die rohen holomorphen Fuchskoordinaten werden nicht als orthonormal erklärt.
Die drei Cartangewichte der tatsächlichen C-Ströme,
$(-1,1,1)/2$, $(1,-1,1)/2$, $(1,1,-1)/2$, prüfen insbesondere den
Faktor $t/2$ unabhängig von einer Endpunktmatrix.

\paragraph{Zustand und Massenkoeffizient besitzen entgegengesetzte Orientierung.}
Für einen symmetrischen Quellenkoeffizienten $X_0$ und eine duale symmetrische
Massenform $K_0$ lautet der faserweise Transport
\begin{equation}\label{eq:massdeterminantcarry}
 \begin{aligned}
 X(z)&=S_C(z)X_0S_C(z)^{\mathsf T},
 &\det X(z)&=d(z)\det X_0,\\
 K(z)&=S_C(z)^{-\mathsf T}K_0S_C(z)^{-1},
 &\det K(z)&=d(z)^{-1}\det K_0.
 \end{aligned}
\end{equation}
Denn $\det S_C=d^{1/2}$. Die beiden Vorzeichen sind daher keine
wählbaren Konventionen innerhalb derselben Zuordnung: Die Paarzustände und
ihre kontrahierenden Koeffizienten sind dual. Das skalare Matrixelement
$v^{\mathsf T}K v$ bleibt erhalten, wenn gleichzeitig $v\mapsto S_Cv$.

Für $\Psi(0)=I$ ist $d=1-z^4$. Auf dem ursprünglichen positiv orientierten
Kreis um $+1$ trägt die Determinantenlinie des invertierbaren Paarzustands
somit Windung $+1$, diejenige der dualen Massenform Windung $-1$.
Bei umgekehrter Kontur kehren sich beide Vorzeichen um. Eine konstante
Phasenänderung des ursprünglichen Basisrahmens ändert diese Windungen nicht.
Eine über der Punktur selbst singuläre, windende Rahmentransformation würde
dagegen die erklärte logarithmische Erweiterung ändern und ist hier kein
unbemerkter Basiswechsel.

Der ursprüngliche schwere Tensor ist für $\epsilon>0$ invertierbar:
\[
 K_0=M_R/M_3=\operatorname{diag}(\epsilon/3,2\epsilon/3,1),
 \qquad \det K_0=2\epsilon^2/9.
\]
Seine bereits berechnete native Einsetzung hat dieselbe nichtverschwindende
Determinante bis auf eine konstante Basisphase. Die Identität gilt für jede
korrekt basierte invertierbare Einsetzung und ersetzt nicht die noch zu
begründende physische Basiszuordnung. Der leichte Seesaw-Ausgang hat dagegen
Rang zwei und Determinante null; für ihn wird keine solche
Determinantenwindung behauptet.

\textbf{Zwei trennende Kontrollen.} Würde man den Faktor $d^{-1/2}$ im
tatsächlichen Außenpotenz-Lift weglassen, entstünde für $\det X$ Windung
vier. Würde man die C-Felder direkt mit dem ursprünglichen Vektordreier
$\Psi$ transportieren, entstünde Windung zwei. Der korrekte Spintransport
liefert Windung eins. Damit trifft die primitive Klasse die vollständige
Tensorwirkung und nicht nur eine Zahl am Endpunkt.

Die \emph{Determinantenlinie} schließt nach einem Punkturumlauf. Ein
allgemeiner voller Tensor wird dabei durch die Familienmonodromie verändert;
der unabhängige direkte Wegtest erhält ihn erst nach drei Umläufen zurück.
Die Gleichung setzt also keinen unveränderten Massentensor nach jeder
einzelnen Runde voraus. Ebenso wenig beweist sie schon einen
vierdimensionalen Vortex-Indexsatz oder die Gleichheit des vollständigen
P1-Nahtoperators mit diesem Familienoperator.

\paragraph{Die sechs Paarkanäle behalten die vollständige Verbindung.}
Die induzierte Verbindung auf $\operatorname{Sym}^2\mathbb C^3$ ist
\begin{equation}\label{eq:sixconnection}
 \mathcal L(X)=BX+XB^{\mathsf T}.
\end{equation}
Aus ihrer Wirkung auf $E_{ii}$ liest man $2B_{ii}$ auf dem diagonalen
Eintrag und $B_{ki}$ auf dem Eintrag $(k,i)$ ab. Damit sind sämtliche neun
komplexen Einträge von $B$ bekannt. Es folgt
\[
 \boxed{A=(\tr B)I-JB^{\mathsf T}J^{\mathsf T}.}
\]
Die Abbildung der Verbindung ist daher injektiv. Eine beliebige
$6\times6$-Matrix liegt nicht notwendig in ihrem neundimensionalen Bild;
die Implementierung prüft die vollständige induzierte Form beim Rücklesen.

Für beliebige, auch nichtkommutierende $3\times3$-Matrizen $A,D$ gilt
\begin{equation}\label{eq:pairtraces}
 \Tr_6\mathcal L_A=2\tr A,\qquad
 \Tr_6(\mathcal L_A\mathcal L_D)=5\tr(AD)-(\tr A)(\tr D).
\end{equation}
Dies folgt aus
$\Tr_{\mathrm{Sym}^2}(B\otimes I+I\otimes B)=4\tr B$
und der entsprechenden bilinearen Spuridentität
$5\tr(BC)+\tr B\tr C$. Die Rechnung gilt für ganze Operatoren;
ein Vergleich ihrer Eigenwertmengen wäre dafür nicht ausreichend.

\paragraph{Die neutrale Antwort ist eine Kontraktion genau dieser Verbindung.}
Im festgehaltenen ursprünglichen Fuchsrahmen und mit derselben symmetrischen
RHP-Normalordnung wie in Abschnitt~\ref{sec:fullflavorRR} folgt exakt
\begin{align}
 \langle R(z)\rangle_F
 &=\frac25\left[\Tr_6\mathcal L(z)^2-(\Tr_6\mathcal L(z))^2\right],\label{eq:massRconnection}\\
 \langle R(z)R(w)\rangle_{F,c}
 &=\frac{48}{(z-w)^4}
 +\frac{16}{5(z-w)^2}\left[\Tr_6\bigl(\mathcal L(z)\mathcal L(w)\bigr)
 +\frac32\Tr_6\mathcal L(z)\Tr_6\mathcal L(w)\right].\label{eq:massRRconnection}
\end{align}
Die rechte Seite stimmt mit dem bereits berechneten vollständigen
rationalen Flavor-Defektkern überein. Es wird keine zweite neutrale Quelle
eingesetzt. Die Spuren einer Verbindungsform sind bei beliebigen lokalen
Rahmenwechseln nicht für sich invariant; deren Ward-/Normalordnungsterme
müssen dann ebenfalls transformiert werden. Die Gleichungen gelten im
erklärten gemeinsamen Originalrahmen, nicht als neue eichinvariante
vierdimensionale Wirkung.

\paragraph{Auch der nötige lokale Ableitungsterm lässt sich aus derselben Quelle ergänzen.}
Eine faserweise Darstellung allein berechtigt noch nicht dazu, gewöhnliche
Ableitungen von ortsabhängig gedrehten Feldern einzusetzen. Hier ist die
Korrektur konkret. In einer Spaltenkonvention für die Feldkomponenten sei
\[
 C(z+s)C(z)^{\mathsf T}=s\mathcal V+s^2\mathcal W+\cdots,
 \quad \mathcal V^{\mathsf T}=-\mathcal V,
 \quad \mathcal B_{\rm raw}=\tfrac12(\mathcal W+\mathcal W^{\mathsf T}).
\]
Der konstante Term fehlt wegen der tatsächlichen Nullrelation
$:C_iC_j:=0$. $\mathcal V$ ist der bereits vorhandene antisymmetrische
Grad-drei-Kanal. Für eine Feldverbindung $b$ liefert das Zurücktransportieren
des ersten Felds mit $I-sb+O(s^2)$ vor dem Zusammenführen
\begin{equation}\label{eq:covariantpairjet}
 \boxed{\mathcal B_{\rm cov}
 =\mathcal B_{\rm raw}-\tfrac12(b\mathcal V-\mathcal V b^{\mathsf T}).}
\end{equation}
Bei $C'=GC$ und $b'=G'G^{-1}+GbG^{-1}$ entsteht im rohen symmetrischen
Koeffizienten der Zusatz
$\tfrac12(G'\mathcal V G^{\mathsf T}-G\mathcal V G'^{\mathsf T})$.
Gleichung~\eqref{eq:covariantpairjet} hebt ihn exakt auf:
$\mathcal B'_{\rm cov}=G\mathcal B_{\rm cov}G^{\mathsf T}$.
Die Spaltenkonvention der Feldkomponenten ist dabei von den dualen
Koeffizienten zur Kombination von Zuständen zu unterscheiden.
Wegen des verschwindenden konstanten OPE-Terms wird an dieser Stelle kein
zweiter Verbindungsjet benötigt. Im festen gewöhnlichen Ableitungsrahmen
bleibt der Grad-drei-Beitrag aber real vorhanden. Eine reine Kürzung auf
sechs Grad-vier-PBW-Zustände wäre somit nicht unter lokalen Rahmenwechseln
abgeschlossen.

\begin{center}
\begin{tikzpicture}[node distance=5mm]
\node[box,text width=12.6cm] (a) {\textbf{Eine ursprüngliche Verbindung $A$}\\
Vier Marken, vollständige Residuen, festgehaltene Erweiterung};
\node[box,below=of a,text width=12.6cm] (b) {\textbf{Ihr Spintransport und die sechs Paarkanäle}\\
$B=\tfrac12(\tr A)I-J^{\mathsf T}A^{\mathsf T}J$,
\quad $\mathcal L(X)=BX+XB^{\mathsf T}$};
\node[box,below=of b,text width=12.6cm] (c) {\textbf{Drei verbundene Antworten}\\
Primitive Windung \quad vollständiger Tensortransport \quad neutrale $R$-Antwort};
\node[box,below=of c,text width=12.6cm] (d) {\textbf{Physische Auswahl noch herzuleiten}\\
Ursprüngliche Rand-/Higgsbedingung für $X_0$ und ursprüngliche Nahtkopplung};
\draw[arr] (a)--(b);\draw[arr] (b)--(c);\draw[dasharr] (c)--(d);
\end{tikzpicture}
\end{center}

\picturetext{Ein gemeinsames Getriebe bewegt jetzt nachweislich die
Orientierung der Massentabelle und die neutrale Anzeige. Sein mitgeführtes
Vorzeichen sorgt dafür, dass aus einer ursprünglichen Runde auch genau
eine Runde der Paar-Determinante wird. Welche Massentabelle in dieses
Getriebe eingesetzt wird, ist eine andere Frage: Die Transportregel
bewegt jede Anfangstabelle korrekt. Sie wählt deren Werte noch nicht aus.}

\textbf{Bedeutung für die fundamentale Auswahl.}
Die Verbindung zwischen P1-artiger Windung, ursprünglichem Familienweg,
symmetrischem Massensektor und neutraler Antwort ist damit als gemeinsames
internes Transportgesetz geschlossen. Offen bleibt die physische
Ursprungsabbildung: Die homogene Gleichung
$X'=BX+XB^{\mathsf T}$ besitzt für \emph{jede} Anfangsmatrix eine Lösung.
Sie kann daher weder die schwere Textur noch deren Skala auswählen.
Auch rohe Änderungen der Matrixsingulärwerte durch einen nichtunitären
Rahmen sind keine erzeugte Massenhierarchie: Die gleichzeitig transportierte
positive Metrik erhält die normierten Größen.
Gesucht ist nun die ursprüngliche Randbedingung beziehungsweise die
tatsächlich erzeugte geladene Higgsantwort in diesem Kanal. Im letzteren
Fall müsste ihre Quelle in
$X'-BX-XB^{\mathsf T}=\mathcal J_H$
aus der ursprünglichen Wirkung berechnet werden; $\mathcal J_H$ wird hier
nicht zur gewünschten Antwort passend definiert. Dieselbe Wirkung muss
die weiterhin bedingte Nahtkopplung $c_3^2|R\rangle\langle p_0|$ und die
physische Zeitzuordnung tragen.

\file{tfpt_explorer/source_mass_transport.py} implementiert die induzierte
Verbindung, ihre vollständige Umkehrung, den endlichen Transport, die
kovariante lokale Ergänzung und die gemeinsamen Antwortformeln.
Die direkte ursprüngliche $3\times3$-Wegintegration und eine separat
integrierte $6\times6$-Gleichung stimmen auf beiden Umlauforientierungen
mit einem relativen Fehler unter $4\cdot10^{-12}$ überein.
Diese numerische Kontrolle ergänzt die ausgeschriebenen algebraischen
Identitäten; sie ist kein physischer Herkunfts- oder Kondensationsbeweis.

